\begin{remark}[判别式诱导的伴随二次域：递推证书的二次分辨率]\label{rem:pom-discriminant-quadratic-field-tower}
对每个整数 \(q\ge 2\)，令 \(P_q(\lambda)\in\ZZ[\lambda]\) 表示 \(m\mapsto S_q(m)\) 的最小整系数递推特征多项式。
记
\[
\Delta_q^{\mathrm{disc}}:=\mathrm{sqf}\!\bigl(\mathrm{Disc}(P_q)\bigr)\in\ZZ
\]
为其判别式的平方自由部分（含符号），并定义伴随二次域
\[
K_q^{\mathrm{disc}}:=\QQ\bigl(\sqrt{\Delta_q^{\mathrm{disc}}}\bigr).
\]
在低阶可显式实现的三例 \(q=2,3,4\) 中，递推特征多项式分别为
\[
P_2(\lambda)=\lambda^3-2\lambda^2-2\lambda+2,\qquad
P_3(\lambda)=\lambda^3-2\lambda^2-4\lambda+2,\qquad
P_4(\lambda)=\lambda^5-2\lambda^4-7\lambda^3-2\lambda+2,
\]
从而其二次分辨率指纹可直接读出为
\[
\Delta_2^{\mathrm{disc}}=37,\qquad
\Delta_3^{\mathrm{disc}}=3\cdot 47,\qquad
\Delta_4^{\mathrm{disc}}=-985219
\]
（其中 \(q=4\) 的判别式由命题 \ref{prop:pom-a4-galois-s5} 给出）。
因此四阶碰撞核在判别式层面诱导出一条全新的虚二次方向
\(
K_4^{\mathrm{disc}}=\QQ(\sqrt{-985219})
\)，
并与 Newman 边权阈值证书的判别式二次子域 \(K_2=\QQ(\sqrt{-1151})\)（命题 \ref{prop:pom-a4t-newman-k2-class-number}）在算术意义上严格分离。
该伴随二次域序列为“最小递推不仅决定实谱增长常数，还在判别式层面固定二次子方向”提供了统一接口；其局部指纹可由 \(P_q\) 在有限域上的分解型（Frobenius 置换类型）直接投射并审计（参见命题 \ref{prop:pom-a4-galois-s5} 的模 \(p\) 分解证书链）。
\end{remark}

\begin{proposition}[二次分辨率的 Frobenius 奇偶筛：\texorpdfstring{$\left(\frac{37}{p}\right)$}{(37/p)} 决定 \texorpdfstring{$P_2\bmod p$}{P2 mod p} 的奇偶分解型]\label{prop:pom-p2-frobenius-parity-sieve}
\leanverified{paper\_pom\_p2\_frobenius\_parity\_sieve}
令 \(P_2(\lambda)=\lambda^3-2\lambda^2-2\lambda+2\)，并令 \(L_2\) 为其分裂域。
由于 \(P_2\) 在 \(\QQ\) 上不可约且 \(\Delta_2^{\mathrm{disc}}=37\) 非平方，故
\(\mathrm{Gal}(L_2/\QQ)\cong S_3\)，并且其唯一二次子域恰为
\(
K_2^{\mathrm{disc}}=\QQ(\sqrt{37})
\)。
对任意不分歧素数 \(p\nmid 2\cdot 37\)，记 \(\mathrm{Frob}_p\in \mathrm{Gal}(L_2/\QQ)\) 为 Frobenius 元。
则其置换奇偶由 Legendre 符号 \(\left(\frac{37}{p}\right)\) 完全决定，从而 \(P_2(\lambda)\bmod p\) 的分解型满足：
\begin{enumerate}
  \item 若 \(\left(\frac{37}{p}\right)=-1\)，则 \(\mathrm{Frob}_p\) 为奇置换（在 \(S_3\) 中即换位），故 \(P_2\bmod p\) 必为分解型 \((1)(2)\)；
  \item 若 \(\left(\frac{37}{p}\right)=+1\)，则 \(\mathrm{Frob}_p\) 为偶置换，故 \(P_2\bmod p\) 的分解型只能为 \((3)\) 或 \((1)(1)(1)\)。
\end{enumerate}
并且在 Chebotarev 等分布下，条件于 \(\left(\frac{37}{p}\right)=+1\) 的素数子集内，
\((1)(1)(1)\) 与 \((3)\) 的相对频率分别为 \(1/3\) 与 \(2/3\)。
\end{proposition}
\begin{proof}
对 \(S_3\)-三次分裂域，符号 \(\mathrm{sgn}:\mathrm{Gal}(L_2/\QQ)\to\{\pm 1\}\) 的核对应于唯一二次子域 \(K_2^{\mathrm{disc}}\)。
因此对不分歧素数 \(p\nmid 2\cdot 37\)，有
\(
\mathrm{sgn}(\mathrm{Frob}_p)=\left(\frac{37}{p}\right)
\)，
从而 \(\left(\frac{37}{p}\right)=-1\) 当且仅当 \(\mathrm{Frob}_p\) 属于换位共轭类，对应分解型 \((1)(2)\)；
\(\left(\frac{37}{p}\right)=+1\) 当且仅当 \(\mathrm{Frob}_p\) 属于单位元或 \(3\)-循环共轭类，对应分解型 \((1)(1)(1)\) 或 \((3)\)。
无条件密度由 \(S_3\) 的共轭类大小给出：单位元、换位、\(3\)-循环三类密度分别为 \(1/6,1/2,1/3\)；
将其在偶置换子集（密度 \(1/2\)）内归一化即得到所述条件相对频率。证毕。
\end{proof}

\begin{proposition}[\texorpdfstring{$K_4^{\mathrm{disc}}=\QQ(\sqrt{-985219})$}{K4disc} 的类数证书：\(h=195\)]\label{prop:pom-a4-disc-quadratic-class-number}
\leanverified{paper\_pom\_a4\_disc\_quadratic\_class\_number}
令 \(K_4^{\mathrm{disc}}=\QQ(\sqrt{-985219})\)。
则其类数满足
\[
h(K_4^{\mathrm{disc}})=195=3\cdot 5\cdot 13.
\]
\end{proposition}
\begin{proof}
由脚本 \texttt{scripts/exp\_collision\_kernel\_A4\_charpoly\_quadratic\_field\_class\_number.py}，
先对 \(p_7(x)=x^5-2x^4-7x^3-2x+2\) 计算 \(\mathrm{Disc}_x(p_7)\) 并抽取其平方自由部分 \(-985219\)，从而得到二次域 \(K_4^{\mathrm{disc}}\)。
再枚举判别式为 \(-985219\) 的约化本原正定二元二次型，得到其个数为 \(195\)，即为该虚二次域的类数。
相应的整数分解亦由脚本输出给出。证毕。
\end{proof}
